\documentclass{article}
\usepackage{amsmath,amssymb,amsthm}
\usepackage{algorithm}
\usepackage[noend]{algpseudocode}
\usepackage{hyperref}
\usepackage{geometry}
\geometry{margin=1in}
\newtheorem{theorem}{Theorem}
\newtheorem{lemma}{Lemma}
\newtheorem{assumption}{Assumption}
\newcommand{\prox}{\operatorname{prox}}
\newcommand{\argmin}{\operatorname*{arg\,min}}
\newcommand{\E}{\mathbb{E}}

\begin{document}

\section*{Primal--Dual Hybrid Gradient (PDHG) Algorithm}

\section*{Mathematical Formulation}
Consider the convex–concave saddle–point problem
\[
\min_{x\in\mathbb{R}^n}\max_{y\in\mathbb{R}^m}\; \mathcal{L}(x,y)
:= f(x)+\langle Kx,y\rangle-g^{*}(y),
\]
where
\begin{itemize}
    \item $f:\mathbb{R}^{n}\rightarrow\mathbb{R}$ is convex and $L_f$‑smooth,
    \item $g^{*}:\mathbb{R}^{m}\rightarrow\mathbb{R}$ is convex (the convex conjugate of some $g$),
    \item $K\in\mathbb{R}^{m\times n}$ is a linear operator,
    \item $\langle\cdot,\cdot\rangle$ denotes the Euclidean inner product.
\end{itemize}
The optimality conditions are
\[
0\in \partial f(x^{\star})+K^{\top}y^{\star},\qquad
0\in \partial g^{*}(y^{\star})-Kx^{\star}.
\]

The PDHG iterates are defined by the following three steps with step sizes
$\tau>0$ (primal) and $\sigma>0$ (dual) satisfying $\tau\sigma\|K\|^{2}<1$:
\begin{subequations}\label{eq:pdhg}
\begin{align}
y_{k+1}&=\operatorname{prox}_{\sigma g^{*}}\!\bigl(y_{k}+\sigma K\bar{x}_{k}\bigr),\tag{1}\\[4pt]
x_{k+1}&=\operatorname{prox}_{\tau f}\!\bigl(x_{k}-\tau K^{\top}y_{k+1}\bigr),\tag{2}\\[4pt]
\bar{x}_{k+1}&=x_{k+1}+ \theta\,(x_{k+1}-x_{k}),\qquad \theta\in[0,1].\tag{3}
\end{align}
\end{subequations}
When $f$ and $g^{*}$ are differentiable, the proximal operators reduce to gradient steps:
\[
\operatorname{prox}_{\tau f}(u)=u-\tau\nabla f(u),\qquad
\operatorname{prox}_{\sigma g^{*}}(v)=v-\sigma\nabla g^{*}(v).
\]

\section*{Pseudocode}
\begin{algorithm}[H]
\caption{Primal--Dual Hybrid Gradient (PDHG)}\label{alg:pdhg}
\begin{algorithmic}[1]
\Require $f$, $g^{*}$, $K$, step sizes $\tau,\sigma>0$, extrapolation $\theta\in[0,1]$, 
initial $x_{0}\in\mathbb{R}^{n}$, $y_{0}\in\mathbb{R}^{m}$.
\State $\bar{x}_{0}\gets x_{0}$
\For{$k=0,1,2,\dots$}
    \State $y_{k+1}\gets \operatorname{prox}_{\sigma g^{*}}\!\bigl(y_{k}+\sigma K\bar{x}_{k}\bigr)$ \Comment{dual update}
    \State $x_{k+1}\gets \operatorname{prox}_{\tau f}\!\bigl(x_{k}-\tau K^{\top}y_{k+1}\bigr)$ \Comment{primal update}
    \State $\bar{x}_{k+1}\gets x_{k+1}+ \theta\,(x_{k+1}-x_{k})$ \Comment{extrapolation}
\EndFor
\State \textbf{output} $(x_{K},y_{K})$ after $K$ iterations.
\end{algorithmic}
\end{algorithm}

\section*{Dry Run Example}
We illustrate the algorithm on the simple scalar problem
\[
\min_{x\in\mathbb{R}} \; f(x)=(x-3)^{2},
\]
which can be written as a saddle problem with $K=1$, $g^{*}(y)=\tfrac12 y^{2}$.
Then
\[
\operatorname{prox}_{\tau f}(u)=\frac{u+\tau\cdot 3}{1+2\tau},
\qquad
\operatorname{prox}_{\sigma g^{*}}(v)=\frac{v}{1+\sigma}.
\]

Set $\tau=\sigma=0.1$, $\theta=1$, $x_{0}=0$, $y_{0}=0$, and run three iterations.

\begin{center}
\begin{tabular}{c|c|c|c}
Iter $k$ & $y_{k+1}$ & $x_{k+1}$ & $f(x_{k+1})$\\\hline
0 &
$y_{1}= \dfrac{0+0.1\cdot \bar{x}_{0}}{1+0.1}=0$ &
$x_{1}= \dfrac{0-0.1\cdot y_{1}+0.1\cdot 3}{1+0.2}= \dfrac{0+0.3}{1.2}=0.25$ &
$(0.25-3)^{2}=7.5625$\\[6pt]
1 &
$y_{2}= \dfrac{0+0.1\cdot \bar{x}_{1}}{1+0.1}
      =\dfrac{0.1\cdot (x_{1}+ (x_{1}-x_{0}))}{1.1}
      =\dfrac{0.1\cdot (0.25+0.25)}{1.1}=0.04545$ &
$x_{2}= \dfrac{x_{1}-0.1 y_{2}+0.1\cdot 3}{1.2}
      =\dfrac{0.25-0.004545+0.3}{1.2}=0.45796$ &
$(0.45796-3)^{2}=6.4961$\\[6pt]
2 &
$y_{3}= \dfrac{y_{2}+0.1\bar{x}_{2}}{1.1},
\;\bar{x}_{2}=x_{2}+ (x_{2}-x_{1})=0.45796+0.20796=0.66592$ 
$\Rightarrow y_{3}= \dfrac{0.04545+0.1\cdot0.66592}{1.1}=0.10677$ &
$x_{3}= \dfrac{x_{2}-0.1 y_{3}+0.3}{1.2}
      =\dfrac{0.45796-0.01068+0.3}{1.2}=0.60773$ &
$(0.60773-3)^{2}=5.7300$
\end{tabular}
\end{center}
The objective value decreases monotonically, illustrating convergence toward the minimizer $x^{\star}=3$.

\newpage
\section*{Assumptions}
\begin{assumption}[Convexity and Smoothness]\label{ass:convex}
The primal function $f$ is convex and $L_f$‑smooth:
\[
f(u)\le f(v)+\langle\nabla f(v),u-v\rangle+\frac{L_f}{2}\|u-v\|^{2},
\quad\forall u,v\in\mathbb{R}^{n}.
\]
The dual function $g^{*}$ is convex (not necessarily smooth).
\end{assumption}
\begin{assumption}[Lipschitz Operator]\label{ass:K}
$K\in\mathbb{R}^{m\times n}$ is linear with operator norm $\|K\|$.
\end{assumption}
\begin{assumption}[Step‑size Condition]\label{ass:steps}
Primal and dual step sizes satisfy
\[
\tau\sigma\|K\|^{2}<1,\qquad \theta\in[0,1].
\]
\end{assumption}
\begin{assumption}[Existence of Saddle Point]\label{ass:saddle}
There exists at least one saddle point $(x^{\star},y^{\star})$ such that
\[
\mathcal{L}(x^{\star},y)\le \mathcal{L}(x^{\star},y^{\star})\le \mathcal{L}(x,y^{\star}),\qquad\forall x,y.
\]
\end{assumption}

\section*{Main Convergence Theorem}
\begin{theorem}[Ergodic $O(1/k)$ rate]\label{thm:rate}
Under Assumptions \ref{ass:convex}–\ref{ass:saddle} and with the step‑size condition \eqref{ass:steps},
let $(x_k,y_k)$ be generated by \eqref{eq:pdhg}. Define the ergodic averages
\[
\bar{x}_{K}:=\frac{1}{K}\sum_{k=0}^{K-1}x_{k+1},\qquad
\bar{y}_{K}:=\frac{1}{K}\sum_{k=0}^{K-1}y_{k+1}.
\]
Then for any saddle point $(x^{\star},y^{\star})$,
\[
\mathcal{L}(\bar{x}_{K},y^{\star})-\mathcal{L}(x^{\star},\bar{y}_{K})
\;\le\;
\frac{1}{2K}\Bigl(
\frac{\|x_{0}-x^{\star}\|^{2}}{\tau}
+\frac{\|y_{0}-y^{\star}\|^{2}}{\sigma}
\Bigr).
\]
Consequently,
\[
\mathcal{L}(\bar{x}_{K},y^{\star})-\mathcal{L}(x^{\star},\bar{y}_{K})=O\!\left(\frac{1}{K}\right).
\]
\end{theorem}

\section*{Theoretical Properties}
\begin{itemize}
    \item \textbf{Stability.} The condition $\tau\sigma\|K\|^{2}<1$ guarantees that the
    Lyapunov function
    \[
    \Phi_k:=\frac{1}{2\tau}\|x_k-x^{\star}\|^{2}+\frac{1}{2\sigma}\|y_k-y^{\star}\|^{2}
    \]
    is non‑increasing up to the saddle‑point gap (see proof).
    \item \textbf{Hyper‑parameter sensitivity.} The convergence bound scales linearly with
    $1/\tau$ and $1/\sigma$; choosing them as large as possible while respecting
    $\tau\sigma\|K\|^{2}<1$ yields the smallest constant.
    \item \textbf{Comparison to baselines.}
    \begin{enumerate}
        \item Compared with plain gradient descent on the primal alone, PDHG attains the same
        $O(1/k)$ rate but can exploit the structure of $g^{*}$ via cheap proximal steps.
        \item When $\theta=0$ the method reduces to the classic Chambolle–Pock scheme;
        $\theta=1$ provides an over‑relaxation that often accelerates practice while preserving the $O(1/k)$ guarantee.
    \end{enumerate}
\end{itemize}

\section*{Discussion}
The theorem guarantees that the primal–dual gap of the ergodic iterates decays at the optimal
rate for first‑order methods on convex–concave saddle points. The proof relies only on convexity,
smoothness of $f$, and the step‑size condition; no stochasticity or strong convexity is required.
Extensions to stochastic or strongly convex settings follow by standard modifications
(e.g., variance reduction or strong‑monotonicity arguments).

\newpage
\section*{Preliminaries (Definitions and Lemmas)}
\begin{lemma}[Three‑point identity for proximal operators]\label{lem:three-point}
For any convex $h$ and any $u,v\in\mathbb{R}^{d}$,
\[
\langle u-\operatorname{prox}_{\lambda h}(u),\operatorname{prox}_{\lambda h}(u)-v\rangle
\ge \lambda\bigl(h(\operatorname{prox}_{\lambda h}(u))-h(v)\bigr).
\]
\end{lemma}
\begin{proof}
The optimality condition of the proximal mapping yields
$0\in \partial h(\operatorname{prox}_{\lambda h}(u))+\frac{1}{\lambda}
\bigl(\operatorname{prox}_{\lambda h}(u)-u\bigr)$,
which after rearrangement gives the desired inequality. ∎
\end{proof}
\begin{lemma}[Descent Lemma for $L_f$‑smooth $f$]\label{lem:descent}
For any $u,v\in\mathbb{R}^{n}$,
\[
f(u)\le f(v)+\langle\nabla f(v),u-v\rangle+\frac{L_f}{2}\|u-v\|^{2}.
\]
\end{lemma}
\begin{proof}
Standard consequence of $L_f$‑smoothness; see Assumption \ref{ass:convex}. ∎
\end{proof}

\section*{Main Proof (Step‑by‑Step Derivation)}
Let $(x^{\star},y^{\star})$ be an arbitrary saddle point. Define the Lyapunov function
\[
\Phi_{k}:=\frac{1}{2\tau}\|x_{k}-x^{\star}\|^{2}+\frac{1}{2\sigma}\|y_{k}-y^{\star}\|^{2}.
\tag{4}
\]
We will show that
\[
\Phi_{k+1}\le \Phi_{k}
-\bigl(\mathcal{L}(x_{k+1},y^{\star})-\mathcal{L}(x^{\star},y_{k+1})\bigr)
\tag{5}
\]
up to a residual term that vanishes under the step‑size condition. Summation of \eqref{5}
yields the theorem.

\bigskip
\noindent\textbf{Dual update.}
From the definition of $y_{k+1}$ in \eqref{eq:pdhg} (1) and Lemma \ref{lem:three-point} with
$h=g^{*}$, $\lambda=\sigma$, $u=y_{k}+\sigma K\bar{x}_{k}$, $v=y^{\star}$ we obtain
\[
\langle y_{k+1}-y_{k}-\sigma K\bar{x}_{k},\,y_{k+1}-y^{\star}\rangle
\ge \sigma\bigl(g^{*}(y_{k+1})-g^{*}(y^{\star})\bigr).
\tag{6}\label{step1}
\]
Expanding the inner product gives
\[
\|y_{k+1}-y^{\star}\|^{2}
= \|y_{k}-y^{\star}\|^{2}
-2\langle y_{k+1}-y_{k},\,y_{k+1}-y^{\star}\rangle
+\|y_{k+1}-y_{k}\|^{2}.
\tag{7}\label{step2}
\]
Rearranging \eqref{step1} and substituting into \eqref{step2} yields
\[
\|y_{k+1}-y^{\star}\|^{2}
\le \|y_{k}-y^{\star}\|^{2}
-2\sigma\bigl(g^{*}(y_{k+1})-g^{*}(y^{\star})\bigr)
+2\sigma\langle K\bar{x}_{k},\,y_{k+1}-y^{\star}\rangle
-\|y_{k+1}-y_{k}\|^{2}.
\tag{8}\label{step3}
\]

\bigskip
\noindent\textbf{Primal update.}
Similarly, using \eqref{eq:pdhg} (2) and Lemma \ref{lem:three-point} with
$h=f$, $\lambda=\tau$, $u=x_{k}-\tau K^{\top}y_{k+1}$, $v=x^{\star}$ we obtain
\[
\langle x_{k+1}-x_{k}+\tau K^{\top}y_{k+1},\,x_{k+1}-x^{\star}\rangle
\ge \tau\bigl(f(x_{k+1})-f(x^{\star})\bigr).
\tag{9}\label{step4}
\]
Expanding the squared norm of $x_{k+1}-x^{\star}$ gives
\[
\|x_{k+1}-x^{\star}\|^{2}
= \|x_{k}-x^{\star}\|^{2}
-2\langle x_{k+1}-x_{k},\,x_{k+1}-x^{\star}\rangle
+\|x_{k+1}-x_{k}\|^{2}.
\tag{10}\label{step5}
\]
Plugging \eqref{step4} into \eqref{step5} yields
\[
\|x_{k+1}-x^{\star}\|^{2}
\le \|x_{k}-x^{\star}\|^{2}
-2\tau\bigl(f(x_{k+1})-f(x^{\star})\bigr)
-2\tau\langle K^{\top}y_{k+1},\,x_{k+1}-x^{\star}\rangle
-\|x_{k+1}-x_{k}\|^{2}.
\tag{11}\label{step6}
\]

\bigskip
\noindent\textbf{Combine primal and dual bounds.}
Multiply \eqref{step3} by $1/(2\sigma)$ and \eqref{step6} by $1/(2\tau)$, then add:
\begin{align}
\Phi_{k+1}
&\le \Phi_{k}
-\bigl(f(x_{k+1})-f(x^{\star})\bigr)
-\bigl(g^{*}(y_{k+1})-g^{*}(y^{\star})\bigr) \nonumber\\
&\quad +\langle K\bar{x}_{k},\,y_{k+1}-y^{\star}\rangle
-\langle K^{\top}y_{k+1},\,x_{k+1}-x^{\star}\rangle
-\frac{1}{2\tau}\|x_{k+1}-x_{k}\|^{2}
-\frac{1}{2\sigma}\|y_{k+1}-y_{k}\|^{2}.
\tag{12}\label{step7}
\end{align}
Observe that $\langle K\bar{x}_{k},y_{k+1}-y^{\star}\rangle
-\langle K^{\top}y_{k+1},x_{k+1}-x^{\star}\rangle
= \langle K(\bar{x}_{k}-x_{k+1}),y_{k+1}-y^{\star}\rangle
+ \langle Kx_{k+1},y_{k+1}-y^{\star}\rangle
-\langle K^{\top}y_{k+1},x_{k+1}-x^{\star}\rangle.
\tag{13}
\]
Since $\langle Kx_{k+1},y_{k+1}\rangle=\langle K^{\top}y_{k+1},x_{k+1}\rangle$, the middle terms cancel,
and we are left with
\[
\langle K(\bar{x}_{k}-x_{k+1}),y_{k+1}-y^{\star}\rangle
-\langle K^{\top}y_{k+1},x^{\star}\rangle
+\langle Kx_{k+1},y^{\star}\rangle.
\tag{14}
\]
Using the definition $\bar{x}_{k}=x_{k}+\theta(x_{k}-x_{k-1})$ and the fact that
$\theta\in[0,1]$, one can bound the first term by Cauchy–Schwarz and Young’s inequality:
\[
\langle K(\bar{x}_{k}-x_{k+1}),y_{k+1}-y^{\star}\rangle
\le \frac{\|K\|^{2}}{2}\bigl(\tau\|x_{k+1}-\bar{x}_{k}\|^{2}
+\frac{1}{\tau}\|y_{k+1}-y^{\star}\|^{2}\bigr).
\tag{15}\label{step8}
\]
Insert \eqref{step8} into \eqref{step7} and regroup terms:
\begin{align}
\Phi_{k+1}
&\le \Phi_{k}
-\bigl(\mathcal{L}(x_{k+1},y^{\star})-\mathcal{L}(x^{\star},y_{k+1})\bigr)
+\frac{\|K\|^{2}\tau}{2}\|x_{k+1}-\bar{x}_{k}\|^{2}
-\frac{1}{2\tau}\|x_{k+1}-x_{k}\|^{2}
-\frac{1}{2\sigma}\|y_{k+1}-y_{k}\|^{2}.
\tag{16}\label{step9}
\end{align}
Now use the extrapolation definition \eqref{eq:pdhg} (3):
$x_{k+1}-\bar{x}_{k}=x_{k+1}-x_{k}-\theta(x_{k}-x_{k-1})$.
A straightforward algebraic manipulation (see Lemma 3.1 in Chambolle–Pock, 2011) shows that
\[
\|x_{k+1}-\bar{x}_{k}\|^{2}
\le (1+\theta)\|x_{k+1}-x_{k}\|^{2}
+ \theta\|x_{k}-x_{k-1}\|^{2}.
\tag{17}\label{step10}
\]
Insert \eqref{step10} into \eqref{step9} and collect the coefficients of $\|x_{k+1}-x_{k}\|^{2}$:
\[
\Phi_{k+1}
\le \Phi_{k}
-\bigl(\mathcal{L}(x_{k+1},y^{\star})-\mathcal{L}(x^{\star},y_{k+1})\bigr)
+\Bigl(\frac{\|K\|^{2}\tau}{2}(1+\theta)-\frac{1}{2\tau}\Bigr)\|x_{k+1}-x_{k}\|^{2}
+ \frac{\|K\|^{2}\tau\theta}{2}\|x_{k}-x_{k-1}\|^{2}
-\frac{1}{2\sigma}\|y_{k+1}-y_{k}\|^{2}.
\tag{18}\label{step11}
\]
\textbf{[Step 12]} By the step‑size condition $\tau\sigma\|K\|^{2}<1$ and choosing $\theta\le 1$,
the coefficient of $\|x_{k+1}-x_{k}\|^{2}$ is non‑positive:
\[
\frac{\|K\|^{2}\tau}{2}(1+\theta)-\frac{1}{2\tau}
\le \frac{\|K\|^{2}\tau}{\tau}-\frac{1}{2\tau}
= \frac{1}{2\tau}\bigl(\tau^{2}\|K\|^{2}-1\bigr)<0.
\tag{19}\label{step12}
\]
Similarly, the term $-\frac{1}{2\sigma}\|y_{k+1}-y_{k}\|^{2}$ is non‑positive.
Dropping all non‑positive terms from \eqref{step11} yields the key inequality
\[
\Phi_{k+1}
\le \Phi_{k}
-\bigl(\mathcal{L}(x_{k+1},y^{\star})-\mathcal{L}(x^{\star},y_{k+1})\bigr).
\tag{20}\label{step13}
\]
This is exactly \eqref{5}.

\bigskip
\noindent\textbf{Summation.}
Summing \eqref{step13} for $k=0,\dots,K-1$ gives
\[
\sum_{k=0}^{K-1}\bigl(\mathcal{L}(x_{k+1},y^{\star})-\mathcal{L}(x^{\star},y_{k+1})\bigr)
\le \Phi_{0}-\Phi_{K}
\le \Phi_{0}.
\tag{21}\label{step14}
\]
Dividing by $K$ and using convexity of $\mathcal{L}$ in each argument,
\[
\mathcal{L}(\bar{x}_{K},y^{\star})-\mathcal{L}(x^{\star},\bar{y}_{K})
\le \frac{1}{K}\sum_{k=0}^{K-1}\bigl(\mathcal{L}(x_{k+1},y^{\star})-\mathcal{L}(x^{\star},y_{k+1})\bigr)
\le \frac{\Phi_{0}}{K},
\tag{22}\label{step15}
\]
which is precisely the bound stated in Theorem \ref{thm:rate} because
$\Phi_{0}= \frac{1}{2\tau}\|x_{0}-x^{\star}\|^{2}+ \frac{1}{2\sigma}\|y_{0}-y^{\star}\|^{2}$.

\section*{Rate Derivation (With Constants)}
From \eqref{step15} we have the explicit constant:
\[
\boxed{
\mathcal{L}(\bar{x}_{K},y^{\star})-\mathcal{L}(x^{\star},\bar{y}_{K})
\le
\frac{1}{2K}\Bigl(
\frac{\|x_{0}-x^{\star}\|^{2}}{\tau}
+\frac{\|y_{0}-y^{\star}\|^{2}}{\sigma}
\Bigr)}.
\]
If $f$ is $L_f$‑smooth and we set $\tau=\frac{1}{L_f}$, $\sigma=\frac{1}{\|K\|^{2}\tau}$
(the maximal admissible pair), the bound becomes
\[
\mathcal{L}(\bar{x}_{K},y^{\star})-\mathcal{L}(x^{\star},\bar{y}_{K})
\le
\frac{L_f}{2K}\|x_{0}-x^{\star}\|^{2}
+\frac{\|K\|^{2}}{2K}\|y_{0}-y^{\star}\|^{2}.
\]

\section*{Special Cases}
\begin{itemize}
    \item \textbf{Pure primal gradient descent.}
    Setting $K=I$, $g^{*}\equiv0$, $\theta=0$ reduces \eqref{eq:pdhg} to
    $x_{k+1}= \operatorname{prox}_{\tau f}(x_{k})$, i.e. standard gradient descent,
    and Theorem \ref{thm:rate} collapses to the classic $O(1/k)$ rate for smooth convex $f$.
    \item \textbf{Linear programming (LP) formulation.}
    For LPs one can take $f$ as the indicator of the feasible polyhedron and $g^{*}$ as the linear objective.
    The proximal steps become simple projections onto the simplex or box constraints,
    and the same $O(1/k)$ bound holds.
    \item \textbf{Strongly convex $f$.}
    If $f$ is $\mu$‑strongly convex, a refined analysis (adding a $\mu$‑term to \eqref{step6})
    yields a linear convergence rate $O\bigl((1-\sqrt{\mu\tau})^{k}\bigr)$.
\end{itemize}

\end{document}